\ifdefined\gkver\else\def\gkver{student}\fi
\documentclass[\gkver]{gaokaozhenti}
\begin{document}

\chapter{2011年福建理综}
%% paperType: 理综物理部分

\begin{enumerate}
\item
%% number: 13
%% paperName: 2011年福建理综第 13 题 6 分
%% typeId: 1
%% type: 单选题
%% chapter: 曲线运动
%% point: 人造地球卫星
%% method: 无
%% score: 6
%% degree: 650
%% duplicateId: 0
%% body:
“嫦娥二号”是我国月球探测第二期工程的先导星。若测得“嫦娥二号”在月球（可视为密度均匀的球体）表面附近圆形轨道运行的周期 $T$，已知引力常数 $G$，半径为 $R$ 的球体体积公式 $V=\frac{4}{3}\pi R^{3}$，则可估算月球的\xzanswer{A}
\fourchoices[answer=A]
{密度}
{质量}
{半径}
{自转周期}
%% answer: A
%% memo:
\memoanswer{
由万有引力提供向心力有

$G\frac{M_{\text{月}}m}{R^{2}}=m\left(\frac{2\pi}{T}\right)^{2}R$

得月球质量

$M_{\text{月}}=\frac{4\pi^{2}R^{3}}{GT^{2}}$

又 $\rho=\frac{M_{\text{月}}}{V}$，$V=\frac{4}{3}\pi R^{3}$，得月球密度

$\rho=\frac{3\pi}{GT^{2}}$

已知周期 $T$ 和引力常量 $G$，A 正确。
}

\item
%% number: 14
%% paperName: 2011年福建理综第 14 题 6 分
%% typeId: 1
%% type: 单选题
%% chapter: 光学
%% point: 全反射
%% method: 无
%% score: 6
%% degree: 650
%% duplicateId: 0
%% body:
如图，半圆形玻璃砖置于光屏 PQ 的左下方。一束白光沿半径方向从 A 点射入玻璃砖，在 O 点发生反射和折射，折射光在白光屏上呈现七色光带。若入射点由 A 向 B 缓慢移动，并保持白光沿半径方向入射到 O 点，观察到各色光在光屏上陆续消失。在光带未完全消失之前，反射光的强度变化以及光屏上最先消失的光分别是\xzanswer{C}

\onepicture[w=4.5cm]{figs/14.png}

\fourchoices[answer=C]
{减弱，紫光}
{减弱，红光}
{增强，紫光}
{增强，红光}
%% answer: C
%% memo:
\memoanswer{
由于入射点由 A 向 B 缓慢移动，使得入射角不断增大，当入射角大于临界角时发生全反射。而介质对色光的折射率由红光、橙光、黄光…到紫光增大，因而紫光先发生全反射，光屏上最先消失的光是紫光；在光带未完全消失之前，发生全反射的光由紫光到红光逐渐减少，反射光逐渐增强，C 正确。
}

\item
%% number: 15
%% paperName: 2011年福建理综第 15 题 6 分
%% typeId: 1
%% type: 单选题
%% chapter: 交流电
%% point: 交流电
%% method: 无
%% score: 6
%% degree: 650
%% duplicateId: 0
%% body:
图 \ref{2011福建理综15a} 中理想变压器原、副线圈的匝数之比 $n_{1}:n_{2}=5:1$，电阻 $R=20\UO$，$L_{1}$、$L_{2}$ 为规格相同的两只小灯泡，$S_{1}$ 为单刀双掷开关。原线圈接正弦交变电源，输入电压 $u$ 随时间 $t$ 的变化关系如图 \ref{2011福建理综15b} 所示。现将 $S_{1}$ 接 1、$S_{2}$ 闭合，此时 $L_{2}$ 正常发光。下列说法正确的是\xzanswer{D}

\twopicture[numA=甲, numB=乙, labelA=2011福建理综15a, labelB=2011福建理综15b, w=5.2cm]
{figs/15a.png}{figs/15b.png}

\fourchoices[answer=D]
{输入电压 $u$ 的表达式 $u=20\sqrt{2}\sin(50\pi t)\UV$}
{只断开 $S_{1}$ 后，$L_{1}$、$L_{2}$ 均正常发光}
{只断开 $S_{2}$ 后，原线圈的输入功率增大}
{若 $S_{1}$ 换接到 2 后，$R$ 消耗的电功率为 $0.8\UW$}
%% answer: D
%% memo:
\memoanswer{
由题图乙可得周期 $T=0.02\Us$，则 $\omega=\frac{2\pi}{T}=100\pi\Urads$，输入电压表达式 $u=20\sqrt{2}\sin100\pi t\UV$，A 错误；只断开 $S_{1}$，$L_{1}$、$L_{2}$ 断路均不发光，B 错误；只断开 $S_{2}$ 后，负载电阻增大，输出电压不变，输出功率减小，则输入功率减小，C 错误；若 $S_{1}$ 换接到 2 后，加在 $R$ 上电压为 $4\UV$，其消耗的电功率为 $P=U^{2}/R=0.8\UW$，D 正确。
}

\item
%% number: 16
%% paperName: 2011年福建理综第 16 题 6 分
%% typeId: 1
%% type: 单选题
%% chapter: 运动和力
%% point: 传送带
%% method: 无
%% score: 6
%% degree: 450
%% duplicateId: 0
%% body:
如图 \ref{2011福建理综16a} 所示，绷紧的水平传送带始终以恒定速率 $v_{1}$ 运行。初速度大小为 $v_{2}$ 的小物块从与传送带等高的光滑水平地面上的 A 处滑上传送带。若从小物块滑上传送带开始计时，小物块在传送带上运动的 $v-t$ 图像（以地面为参考系）如图 \ref{2011福建理综16b} 所示。已知 $v_{2}>v_{1}$，则\xzanswer{B}

\twopicture[numA=甲, numB=乙, labelA=2011福建理综16a, labelB=2011福建理综16b, w=5.2cm]
{figs/16a.png}{figs/16b.png}

\fourchoices[answer=B]
{$t_{2}$ 时刻，小物块离 A 处的距离达到最大}
{$t_{2}$ 时刻，小物块相对传送带滑动的距离达到最大}
{$0\sim t_{2}$ 时间内，小物块受到的摩擦力方向先向右后向左}
{$0\sim t_{3}$ 时间内，小物块始终受到大小不变的摩擦力作用}
%% answer: B
%% memo:
\memoanswer{
由题图乙可知，在 $t_{1}$ 时刻物块的速度为零，离开 A 点的距离最大，A 错误；$t_{2}$ 时刻，小物块刚好相对传送带静止，此时相对传送带滑动的距离最大，B 正确；$0\sim t_{2}$ 时间内，小物块在摩擦力的作用下先减速再反向加速，摩擦力不变，C 错误；$t_{2}\sim t_{3}$ 时间内，小物块相对传送带静止且随水平传送带一起匀速运动，不受摩擦力作用，D 错误。
}

\item
%% number: 17
%% paperName: 2011年福建理综第 17 题 6 分
%% typeId: 1
%% type: 单选题
%% chapter: 电磁感应
%% point: 电磁感应中的变加速
%% method: 无
%% score: 6
%% degree: 450
%% duplicateId: 0
%% body:
如图，足够长的 U 型光滑金属导轨平面与水平面成 $\theta$ 角（$0<\theta<90^\circ$），其中 MN 平行且间距为 $L$，导轨平面与磁感应强度为 $B$ 的匀强磁场垂直，导轨电阻不计。金属棒 ab 由静止开始沿导轨下滑，并与两导轨始终保持垂直且良好接触，ab 棒接入电路的电阻为 $R$，当流过 ab 棒某一横截面的电量为 $q$ 时，棒的速度大小为 $v$，则金属棒 ab 在这一过程中\xzanswer{B}

\onepicture[w=4.2cm]{figs/17.png}

\fourchoices[answer=B]
{运动的平均速度大小为 $\frac{1}{2}v$}
{下滑位移大小为 $\frac{qR}{BL}$}
{产生的焦耳热为 $qBLv$}
{受到的最大安培力大小为 $\frac{B^{2}L^{2}v}{R}\sin\theta$}
%% answer: B
%% memo:
\memoanswer{
由于 ab 棒下滑切割磁感线，使闭合电路产生感应电流，ab 棒中有电流在磁场中受到安培力作用，棒做加速度减小的加速运动，不是匀加速运动，平均速度不是 $\frac{1}{2}v$，A 错误；

由法拉第电磁感应定律，$E=\frac{\Delta\phi}{\Delta t}$ 及 $\Delta\phi=B\Delta S=BLx$、$\overline{I}=\frac{E}{R}$ 和 $\overline{I}=\frac{q}{\Delta t}$ 得流过 ab 棒某一横截面的电量为 $q$ 时下滑的位移大小

$x=\frac{qR}{BL}$

B 正确；

产生的焦耳热为 $Q=\overline{I}^{2}R\Delta t=\frac{q^{2}R}{\Delta t}$，C 错误；

安培力的大小为 $F=BIL=\frac{B^{2}L^{2}v}{R}$，D 错误。
}

\item
%% number: 18
%% paperName: 2011年福建理综第 18 题 6 分
%% typeId: 1
%% type: 单选题
%% chapter: 运动和力
%% point: 连接体
%% method: 极限
%% score: 6
%% degree: 250
%% duplicateId: 0
%% body:
如图，一不可伸长的轻质细绳跨过滑轮后，两端分别悬挂质量为 $m_{1}$ 和 $m_{2}$ 的物体 A 和 B。若滑轮有一定大小，质量为 $m$ 且分布均匀，滑轮转动时与绳之间无相对滑动，不计滑轮与轴之间的摩擦。设细绳对 A 和 B 的拉力大小分别为 $T_{1}$ 和 $T_{2}$，已知下列四个关于 $T_{1}$ 的表达式中有一个是正确的，请你根据所学的物理知识，通过一定的分析判断正确的表达式是\xzanswer{C}

\onepicture[h=5cm]{figs/18.png}

\fourchoices[answer=C]
{$T_{1}=\frac{(m+2m_{2})m_{1}g}{m+2(m_{1}+m_{2})}$}
{$T_{1}=\frac{(m+2m_{1})m_{2}g}{m+4(m_{1}+m_{2})}$}
{$T_{1}=\frac{(m+4m_{2})m_{1}g}{m+2(m_{1}+m_{2})}$}
{$T_{1}=\frac{(m+4m_{1})m_{2}g}{m+4(m_{1}+m_{2})}$}
%% answer: C
%% memo:
\memoanswer{
利用极限的思维方式，若滑轮的质量 $m=0$，则细绳对 A 和 B 的拉力大小 $T_{1}$ 和 $T_{2}$ 相等为 $T$。假设 $m_{1}>m_{2}$，A 和 B 一起运动的加速度为 $a$，根据牛顿第二定律分别对 A、B 有：$m_{1}g-T=m_{1}a$、$T-m_{2}g=m_{2}a$，联立解得：$T=\frac{2m_{1}m_{2}g}{m_{1}+m_{2}}$，分析判断可知 C 正确。
}

\item
%% number: 19
%% paperName: 2011年福建理综第 19 题 12 分
%% typeId: 4
%% type: 实验题
%% chapter: 电路
%% point: 伏安法测电阻
%% method: 无
%% score: 12
%% degree: 450
%% duplicateId: 0
%% body:
\begin{enumerate}
	\item
	（6 分）某实验小组在利用单摆测定当地重力加速度的实验中：
	\begin{enumerate}
		\item
		用游标卡尺测定摆球的直径，测量结果如图所示，则该摆球的直径为\tkanswer[1.5cm]{0.97}$\Ucm$（0.96、0.98 均可）。
		\item
		小组成员在实验过程中有如下说法，其中正确的是\tkanswer[1cm]{C}。（填选项前的字母）

		（A）把单摆从平衡位置拉开 $30^\circ$ 的摆角，并在释放摆球的同时开始计时

		（B）测量摆球通过最低点 100 次的时间 $t$，则单摆周期为 $\frac{t}{100}$

		（C）用悬线的长度加摆球的直径作为摆长，代入单摆周期公式计算得到的重力加速度值偏大

		（D）选择密度较小的摆球，测得的重力加速度值误差较小
	\end{enumerate}

	\onepicture[w=5cm]{figs/19.png}

	\item
	（12 分）某同学在探究规格为“$6\UV$，$3\UW$”的小电珠伏安特性曲线实验中：
	\begin{enumerate}
		\item
		在小电珠接入电路前，使用多用电表直接测量小电珠的电阻，则应将选择开关旋至\tkanswer[1.2cm]{D}档进行测量。（填选项前的字母）

		（A）直流电压 $10\UV$　　（B）直流电流 $5\UmA$　　（C）欧姆 $\times100$　　（D）欧姆 $\times1$
		\item
		该同学采用图 \ref{2011福建理综19a} 所示的电路进行测量。图中 $R$ 为滑动变阻器（阻值范围 $0\sim20\UO$，额定电流 $1.0\UA$），$L$ 为待测小电珠，V 为电压表（量程 $6\UV$，内阻 $20\UkO$），A 为电流表（量程 $0.6\UA$，内阻 $1\UO$），$E$ 为电源（电动势 $8\UV$，内阻不计），$S$ 为开关。

		Ⅰ．在实验过程中，开关 $S$ 闭合前，滑动变阻器的滑片 $P$ 应置于最\tkanswer[1cm]{左}端；（填“左”或“右”）

		Ⅱ．在实验过程中，已知各元器件均无故障，但闭合开关 $S$ 后，无论如何调节滑片 $P$，电压表和电流表的示数总是调不到零，其原因是\tkanswer[0.8cm]{1}点至\tkanswer[0.8cm]{5}点的导线没有连接好；（图 \ref{2011福建理综19a} 中的黑色小圆点表示接线点，并用数字标记，空格中请填写图甲中的数字，如“2 点至 3 点”的导线）

		Ⅲ．该同学描绘出小电珠的伏安特性曲线示意图如图 \ref{2011福建理综19b} 所示，则小电珠的电阻值随工作电压的增大而\tkanswer[1.2cm]{增大}。（填“不变”、“增大”或“减小”）
	\end{enumerate}

	\twopicture[numA=甲, numB=乙, labelA=2011福建理综19a, labelB=2011福建理综19b, w=5cm, align=right]
	{figs/19a.png}{figs/19b.png}
\end{enumerate}
%% answer: 
\jdanswer{
\begin{enumerate}
	\item
	① $0.97\Ucm$（0.96、0.98 均可）；② C。

	\item
	① D；② Ⅰ．左，Ⅱ．1 点至 5 点（或 5 点至 1 点），Ⅲ．增大。

\end{enumerate}
}
%% memo:
\memoanswer{
（1）①从题图可得摆球的直径为 $9\Umm+7\times0.1\Umm=9.7\Umm=0.97\Ucm$。（由于读图的误差可能读成 $0.96\Ucm$、$0.98\Ucm$）

②利用单摆测定当地重力加速度为了减小误差应该从平衡位置开始计时，A 错误；摆球通过最低点 100 次，完成了 50 次全振动，周期为 $\frac{t}{50}$，B 错误；由 $g=\frac{4\pi^{2}l}{T^{2}}$，用悬线的长度加摆球的直径作为摆长，摆长 $l$ 增大了，计算得到的重力加速度值偏大，C 正确；密度较小的摆球，受到的空气的相对阻力影响大，使测得的重力加速度值误差增大，D 错误。

（2）①由于小电珠工作时的电阻为 $R=U^{2}/P=12\UO$，不工作时的电阻小于 $12\UO$，应该用欧姆挡的“$\times1$”挡。D 正确。

②Ⅰ．为了保护电表，使开始的时候分得电压最小为零，$P$ 应置于最左端。

Ⅱ．闭合开关 $S$ 后，无论如何调节滑片 $P$，电压表和电流表的示数总是调不到零，说明滑动变阻器的分压不起作用，1、5 点断路，即 1 点至 5 点导线没有连接好。

Ⅲ．由 $R=U/I$ 和题图乙可知，$U$ 增大时，$R$ 增大。
}

\item
%% number: 20
%% paperName: 2011年福建理综第 20 题 15 分
%% typeId: 6
%% type: 计算题
%% chapter: 电场
%% point: 带电粒子的往复运动
%% method: 无
%% score: 15
%% degree: 450
%% duplicateId: 0
%% body:
反射式调管是常用的微波器械之一，它利用电子团在电场中的震荡来产生微波，其震荡原理与下述过程类似。如图所示，在虚线 MN 两侧分别存在着方向相反的两个匀强电场，一带电微粒从 A 点由静止开始，在电场力作用下沿直线在 A、B 两点间往返运动。已知电场强度的大小分别是 $E_{1}=2.0\times10^{3}\UNC$ 和 $E_{2}=4.0\times10^{3}\UNC$，方向如图所示，带电微粒质量 $m=1.0\times10^{-20}\Ukg$，带电量 $q=-1.0\times10^{-9}\UC$，A 点距虚线 MN 的距离 $d_{1}=1.0\Ucm$，不计带电微粒的重力，忽略相对论效应。求：
\begin{enumerate}
	\item
	B 点到虚线 MN 的距离 $d_{2}$；
	\item
	带电微粒从 A 点运动到 B 点所经历的时间 $t$。
\end{enumerate}

\onepicture[align=right, w=5cm]{figs/20.png}

%% answer: 
\jdanswer{
\begin{enumerate}
	\item
	$0.50\Ucm$

	\item
	$1.5\times10^{-8}\Us$

\end{enumerate}
}
%% memo:
\memoanswer{
（1）带电粒子由 A 运动到 B 的过程中，由动能定理有

$\left|q\right|E_{1}d_{1}-\left|q\right|E_{2}d_{2}=0$ ①

由①式解得

$d_{2}=\frac{E_{1}}{E_{2}}d_{1}=0.50\Ucm$ ②

（2）设微粒在虚线 MN 两侧的加速度大小分别为 $a_{1}$、$a_{2}$，由牛顿第二定律有

$\left|q\right|E_{1}=ma_{1}$ ③

$\left|q\right|E_{2}=ma_{2}$ ④

设微粒在虚线 MN 两侧运动的时间分别为 $t_{1}$、$t_{2}$，由运动学公式有

$d_{1}=\frac{1}{2}a_{1}t_{1}^{2}$ ⑤

$d_{2}=\frac{1}{2}a_{2}t_{2}^{2}$ ⑥

又

$t=t_{1}+t_{2}$ ⑦

联立②③④⑤⑥⑦式解得

$t=1.5\times10^{-8}\Us$
}

\item
%% number: 21
%% paperName: 2011年福建理综第 21 题 19 分
%% typeId: 6
%% type: 计算题
%% chapter: 功和能
%% point: 动能定理
%% method: 无
%% score: 19
%% degree: 250
%% duplicateId: 0
%% body:
如图为某种鱼饵自动投放器中的投饵管装置示意图，其下半部 AB 是一长为 $2R$ 的竖直细管，上半部 BC 是半径为 $R$ 的四分之一圆弧弯管，管口沿水平方向，AB 管内有一原长为 $R$、下端固定的轻质弹簧。投饵时，每次总将弹簧长度压缩到 $0.5R$ 后锁定，在弹簧上端放置一粒鱼饵，解除锁定，弹簧可将鱼饵弹射出去。设质量为 $m$ 的鱼饵到达管口 C 时，对管壁的作用力恰好为零。不计鱼饵在运动过程中的机械能损失，且锁定和解除锁定时，均不改变弹簧的弹性势能。已知重力加速度为 $g$。求：
\begin{enumerate}
	\item
	质量为 $m$ 的鱼饵到达管口 C 时的速度大小 $v_{1}$；
	\item
	弹簧压缩到 $0.5R$ 时的弹性势能 $E_{p}$；
	\item
	已知地面与水面相距 $1.5R$，若使该投饵管绕 AB 管的中轴线 $OO'$ 在 $90^\circ$ 角的范围内来回缓慢转动，每次弹射时只放置一粒鱼饵，鱼饵的质量在 $\frac{2}{3}m$ 到 $m$ 之间变化，且均能落到水面。持续投放足够长时间后，鱼饵能够落到水面的最大面积 $S$ 是多少？
\end{enumerate}

\onepicture[align=right, h=5.5cm]{figs/21.png}

%% answer: 
\jdanswer{
\begin{enumerate}
	\item
	$\sqrt{gR}$

	\item
	$3mgR$

	\item
	$8.25\pi R^{2}$

\end{enumerate}
}
%% memo:
\memoanswer{
（1）质量为 $m$ 的鱼饵到达管口 C 时做圆周运动的向心力完全由重力提供，则

$mg=m\frac{v_{1}^{2}}{R}$ ①

由①式解得

$v_{1}=\sqrt{gR}$ ②

（2）弹簧的弹性势能全部转化为鱼饵的机械能，由机械能守恒定律有

$E_{p}=mg(1.5R+R)+\frac{1}{2}mv_{1}^{2}$ ③

由②③式得

$E_{p}=3mgR$ ④

（3）不考虑因缓慢转动装置对鱼饵速度大小的影响，质量为 $m$ 的鱼饵离开管口 C 后做平抛运动，设经过 $t$ 时间落到水面上，离 $OO'$ 的水平距离为 $x_{1}$，由平抛运动规律有

$4.5R=\frac{1}{2}gt^{2}$ ⑤

$x_{1}=v_{1}t+R$ ⑥

由⑤⑥式解得

$x_{1}=4R$ ⑦

当鱼饵的质量为 $\frac{2}{3}m$ 时，设其到达管口 C 时速度大小为 $v_{2}$，由机械能守恒定律有

$E_{p}=\frac{2}{3}mg(1.5R+R)+\frac{1}{2}\left(\frac{2}{3}m\right)v_{2}^{2}$ ⑧

由④⑧式解得

$v_{2}=2\sqrt{gR}$ ⑨

质量为 $\frac{2}{3}m$ 的鱼饵落到水面上时，设离 $OO'$ 的水平距离为 $x_{2}$，则

$x_{2}=v_{2}t+R$ ⑩

由⑤⑨⑩式解得

$x_{2}=7R$

鱼饵能够落到水面的最大面积 $S$

$S=\frac{1}{4}(\pi x_{2}^{2}-\pi x_{1}^{2})=\frac{33}{4}\pi R^{2}$（或 $8.25\pi R^{2}$）
}

\item
%% number: 22
%% paperName: 2011年福建理综第 22 题 20 分
%% typeId: 6
%% type: 计算题
%% chapter: 磁场
%% point: 带电粒子在复合场中的运动
%% method: 无
%% score: 20
%% degree: 250
%% duplicateId: 0
%% body:
如图 \ref{2011福建理综22a}，在 $x<0$ 的空间中存在沿 $y$ 轴负方向的匀强电场和垂直于 $xoy$ 平面向里的匀强磁场，电场强度大小为 $E$，磁感应强度大小为 $B$。一质量为 $m$、电荷量为 $q$（$q>0$）的粒子从坐标原点 O 处，以初速度 $v_{0}$ 沿 $x$ 轴正方向射入，粒子的运动轨迹见图 \ref{2011福建理综22a}，不计粒子的重力。
\begin{enumerate}
	\item
	求该粒子运动到 $y=h$ 时的速度大小 $v$；
	\item
	现只改变入射粒子初速度的大小，发现初速度大小不同的粒子虽然运动轨迹（$y-x$ 曲线）不同，但具有相同的空间周期性，如图 \ref{2011福建理综22b} 所示；同时，这些粒子在 $y$ 轴方向上的运动（$y-t$ 关系）是简谐运动，且都有相同的周期 $T=\frac{2\pi m}{qB}$。
	\begin{enumerate}
		\item
		求粒子在一个周期 $T$ 内，沿 $x$ 轴方向前进的距离 $S$；
		\item
		当入射粒子的初速度大小为 $v_{0}$ 时，其 $y-t$ 图像如图 \ref{2011福建理综22c} 所示，求该粒子在 $y$ 轴方向上做简谐运动的振幅 $A$，并写出 $y-t$ 的函数表达式。
	\end{enumerate}
\end{enumerate}

\threepicture[numA=甲, numB=乙, numC=丙, labelA=2011福建理综22a, labelB=2011福建理综22b, labelC=2011福建理综22c, w=4.8cm]
{figs/22a.png}{figs/22b.png}{figs/22c.png}

%% answer: 
\jdanswer{
\begin{enumerate}
	\item
	$\sqrt{v_{0}^{2}-\frac{2qEh}{m}}$

	\item
	Ⅰ、$\frac{2\pi mE}{qB^{2}}$；Ⅱ、$A=\frac{m}{qB}\left(v_{0}-\frac{E}{B}\right)$，$y=\frac{m}{qB}\left(v_{0}-\frac{E}{B}\right)\left(1-\cos\frac{qB}{m}t\right)$

\end{enumerate}
}
%% memo:
\memoanswer{
（1）由于洛伦兹力不做功，只有电场力做功，由动能定理有

$-qEh=\frac{1}{2}mv^{2}-\frac{1}{2}mv_{0}^{2}$ ①

由①式解得

$v=\sqrt{v_{0}^{2}-\frac{2qEh}{m}}$ ②

（2）Ⅰ、由图乙可知，所有粒子在一个周期 $T$ 内沿 $x$ 轴方向前进的距离相同，即都等于恰好沿 $x$ 轴方向匀速运动的粒子在 $T$ 时间内前进的距离。设粒子恰好沿 $x$ 轴方向匀速运动的速度大小为 $v_{1}$，则

$qv_{1}B=qE$ ③

又

$S=v_{1}T$ ④

式中 $T=\frac{2\pi m}{qB}$，由③④式解得

$S=\frac{2\pi mE}{qB^{2}}$ ⑤

Ⅱ、设粒子在 $y$ 方向上的最大位移为 $y_{m}$（图丙曲线的最高点处），对应的粒子运动速度大小为 $v_{2}$（方向沿 $x$ 轴），因为粒子在 $y$ 方向上的运动为简谐运动，因而在 $y=0$ 和 $y=y_{m}$ 处粒子所受的合外力大小相等，方向相反，则

$qv_{0}B-qE=-(qv_{2}B-qE)$ ⑥

由动能定理有

$-qEy_{m}=\frac{1}{2}mv_{2}^{2}-\frac{1}{2}mv_{0}^{2}$ ⑦

又

$A=\frac{1}{2}y_{m}$ ⑧

由⑥⑦⑧式解得

$A=\frac{m}{qB}\left(v_{0}-\frac{E}{B}\right)$

可写出图丙曲线满足的简谐运动 $y-t$ 函数表达式为

$y=\frac{m}{qB}\left(v_{0}-\frac{E}{B}\right)\left(1-\cos\frac{qB}{m}t\right)$
}

\item
%% number: 28
%% paperName: 2011年福建理综第 28 题 6 分
%% typeId: 1
%% type: 单选题
%% chapter: 气体性质
%% point: 热力学第一定律
%% method: 无
%% score: 6
%% degree: 650
%% duplicateId: 0
%% body:
【物理——选修 3-3】（本题共有两小题，每小题 6 分，共 12 分。每小题只有一个选项符合题意）

\textbf{（1）}如图，曲线 M、N 分别表示晶体和非晶体在一定压强下的熔化过程，图中横轴表示时间 $t$，纵轴表示温度 $T$。从图中可以确定的是\tkanswer[1cm]{B}。（填选项前的字母）

\onepicture[w=4.5cm]{figs/28.png}

\fourchoices
{晶体和非晶体均存在固定的熔点 $T_{0}$}
{曲线 M 的 bc 段表示固液共存状态}
{曲线 M 的 ab 段、曲线 N 的 ef 段均表示固态}
{曲线 M 的 cd 段、曲线 N 的 fg 段均表示液态}

\textbf{（2）}一定量的理想气体在某一过程中，从外界吸收热量 $2.5\times10^{4}\UJ$，气体对外界做功 $1.0\times10^{4}\UJ$，则该理想气体的\xzanswer{D}。（填选项前的字母）

\fourchoices[answer=D]
{温度降低，密度增大}
{温度降低，密度减小}
{温度升高，密度增大}
{温度升高，密度减小}
%% answer: D
%% memo:
\memoanswer{
（1）由题图 M 曲线表示的晶体具有一定的熔点温度 $T_{0}$，N 曲线表示的非晶体没有一定的熔化温度。A 错误；对于 M 曲线：ab 段为固态温升过程，bc 段为固态、液态共存的熔化过程，cd 段为液态温升过程。B 正确；而 N 曲线 efg 段固液都有，C、D 错误。

（2）由热力学第一定律 $Q+W=\Delta U$，可知 $\Delta U=1.5\times10^{4}\UJ$，内能增加，温度升高；理想气体对外界做功，体积增大，密度减小。D 正确。
}

\item
%% number: 29
%% paperName: 2011年福建理综第 29 题 6 分
%% typeId: 1
%% type: 单选题
%% chapter: 运动和力
%% point: 动量守恒定律
%% method: 无
%% score: 6
%% degree: 650
%% duplicateId: 0
%% body:
【物理——选修 3-5】（本题共有两小题，每小题 6 分，共 12 分。每小题只有一个选项符合题意）

\textbf{（1）}爱因斯坦提出了光量子概念并成功地解释光电效应的规律而获得 1921 年的诺贝尔物理学奖。某种金属逸出光电子的最大初动能 $E_{km}$ 与入射光频率 $\nu$ 的关系如图所示，其中 $\nu_{0}$ 为极限频率。从图中可以确定的是\tkanswer[1cm]{D}。（填选项前的字母）

\onepicture[w=4cm]{figs/29.png}

\fourchoices
{逸出功与 $\nu$ 有关}
{$E_{km}$ 与入射光强度成正比}
{$\nu<\nu_{0}$ 时，会逸出光电子}
{图中直线的斜率与普朗克常量有关}

\textbf{（2）}在光滑水平面上，一质量为 $m$，速度大小为 $v$ 的 A 球与质量为 $2m$ 静止的 B 球碰撞后，A 球的速度方向与碰撞前相反。则碰撞后 B 球的速度大小可能是\xzanswer{A}。（填选项前的字母）

\fourchoices[answer=A]
{$0.6v$}
{$0.4v$}
{$0.3v$}
{$0.2v$}
%% answer: A
%% memo:
\memoanswer{
（1）入射光的频率必须大于金属产生光电效应的极限频率，才会逸出光电子，C 错误；而金属产生光电效应的逸出功是由金属自身的性质决定，与入射光频率 $\nu$ 无关，A 错误；由光电效应方程 $h\nu=w_{0}+E_{km}$，逸出光电子的初动能 $E_{km}$ 决定于入射光的频率 $\nu$，与入射光的强度无关，B 错误；图中直线的斜率 $\frac{E_{km}}{\nu}=h-\frac{w_{0}}{\nu}=h\left(1-\frac{\nu_{0}}{\nu}\right)$，斜率与普朗克常量 $h$ 有关，D 正确。

（2）由动量守恒定律：$mv=mv_{A}+2mv_{B}$，而 $v_{A}<0$，则 $v_{B}>0.5v$，所以 A 正确。
}

\end{enumerate}

\end{document}
